\chapter{Before You Install: Kernel Prerequisites}\label{ch:prereq}
The package can be installed on any \distroname{} host. The \emph{sensor} will only do its full job on a
kernel that exposes BTF and has the BPF Linux Security Module enabled. Check this first; it takes a minute
and saves an hour.

\section{The checklist}
\resetsteps
\step{Confirm BTF is available}
CO-RE (compile once, run everywhere) relies on kernel type information.
\begin{lstlisting}[style=console]
$ ls -l /sys/kernel/btf/vmlinux
-r--r--r-- 1 root root 5321944 ... /sys/kernel/btf/vmlinux
\end{lstlisting}
If the file is missing, the sensor cannot relocate its programs on this kernel.

\step{Confirm the kernel was built with BPF LSM}
\begin{lstlisting}[style=console]
$ zgrep CONFIG_BPF_LSM /proc/config.gz  # or: grep CONFIG_BPF_LSM /boot/config-$(uname -r)
CONFIG_BPF_LSM=y
\end{lstlisting}

\step{Check whether \texttt{bpf} is in the active LSM list}
\begin{lstlisting}[style=console]
$ cat /sys/kernel/security/lsm
lockdown,capability,landlock,yama,apparmor,bpf
\end{lstlisting}
The trailing \texttt{bpf} means the in-kernel file-access block (\texttt{kb\_lsm\_file\_open}) can attach.
Without it, the sensor still loads its tracing programs; only the LSM program fails to attach, and the
sensor's startup log says so rather than silently doing nothing.

\section{Enabling BPF LSM when it is missing}
Add \texttt{bpf} to the kernel's \texttt{lsm=} boot parameter, regenerate the boot configuration, reboot.

\ifarch
\subsection*{Arch Linux (GRUB)}
\begin{lstlisting}[style=console]
$ sudo $EDITOR /etc/default/grub
  GRUB_CMDLINE_LINUX_DEFAULT="quiet lsm=landlock,lockdown,yama,apparmor,bpf"
$ sudo grub-mkconfig -o /boot/grub/grub.cfg
$ sudo reboot
\end{lstlisting}
\begin{kbnote}[SYSTEMD-BOOT / OTHER LOADERS]
If you use \texttt{systemd-boot}, append the same \texttt{lsm=...} string to the \texttt{options} line of
your entry under \file{/boot/loader/entries/}. Arch's own kernels build BPF LSM in; you only need the
parameter. Keep \texttt{apparmor} in the list only if you actually run it; listing a module your kernel
lacks is harmless but pointless.
\end{kbnote}
\else
\subsection*{Ubuntu 24.04 (GRUB)}
\begin{lstlisting}[style=console]
$ sudo $EDITOR /etc/default/grub
  GRUB_CMDLINE_LINUX_DEFAULT="quiet splash lsm=landlock,lockdown,yama,apparmor,bpf"
$ sudo update-grub
$ sudo reboot
\end{lstlisting}
\begin{kbnote}[UBUNTU 24.04 DEMO VM]
The reference Ubuntu VM runs kernel \texttt{6.8.0-139-generic} and already reports
\texttt{lockdown,capability,landlock,yama,apparmor,bpf} with BTF present, so no change is needed there.
\end{kbnote}
\fi

\step{Verify after reboot}
\begin{lstlisting}[style=console]
$ cat /sys/kernel/security/lsm
\end{lstlisting}
\texttt{bpf} must appear at the end of the list.

\section{What runs where}
\begin{center}
\small
\begin{tabularx}{\textwidth}{@{}L{3.6cm}L{2.6cm}Y@{}}
\toprule
\textbf{Component} & \textbf{Privilege} & \textbf{Why}\\
\midrule
\cmd{kbd\_sensor} & root & Loads eBPF, attaches LSM and uprobes; needs CAP\_BPF, CAP\_PERFMON, CAP\_SYS\_ADMIN (exact set is kernel-dependent).\\
\cmd{kbd} & root & Owns \file{/run/kb}, \file{/var/lib/kbd}, binds the sockets.\\
\cmd{kbopd} (service) & unprivileged-style hardening & Only a loopback listener and a child process.\\
\cmd{sshd@kb-operator} & root, then drops to \texttt{operator} & Standard \cmd{sshd} privilege separation.\\
\bottomrule
\end{tabularx}
\end{center}

\begin{kbwarn}[A VM OR CONTAINER SHARES THE HOST KERNEL ASSUMPTIONS]
Inside a container the sensor sees the \emph{host} kernel. Inside a VM it sees the \emph{guest} kernel, which
is what you want for the demo. Whichever you use, run the checklist above \emph{inside} it.
\end{kbwarn}
